% GENERATED FILE. Edit canonical lessons, then run npm run generate:books.
% canonical-source-hash: fnv1a64:56cc2b28fd1d3b56
% canonical-lessons: TA-C61-hunger, TA-C61-thirst, TA-C61-sleep, TA-C61-fever, TA-C61-pain

\chapter{How the Body Feels}
\label{ch:ta-how-the-body-feels}
\hlchaptermodality{61}

\begin{chapteropening}
\textbf{By the end of this chapter:} \emph{I can say I am hungry, thirsty, sleepy, feverish or in pain in Tamil.}

Complete TA-C61-pain and say all five words in order.
\end{chapteropening}

\section[pasi]{\ta{பசி} (pasi) — hunger}

\label{lesson:TA-C61-hunger}



\begin{quote}
Before the new one: what did \emph{tuṭaippam} mean, and what did \emph{sīppu} mean?
\end{quote}

\subsection*{You'll want to know: \ta{பசி}}
\textbf{\ta{பசி}} (\emph{pasi}) — \textquotedblleft{}hunger\textquotedblright{}.

Hunger as a thing that arrives, rather than something you are. Kannada makes it \ka{ಹಸಿವು} (\emph{hasivu}) — the \emph{p} traded for \emph{h} one more time, exactly as it was between \ta{பசு} and \ka{ಹಸು}. Telugu went elsewhere, to \te{ఆకలి} (\emph{ākali}).

Set \ta{பசி} beside \ta{பசு} and only the last mark separates them: one ends on the \ta{ி}, the other on the \ta{ு}. They are not related, and the difference is worth practising until it stops needing thought.

\begin{grammarlens}[title={What you've built}]
The first of five feelings.
\end{grammarlens}

\subsection*{Your turn}
Say these aloud:

\begin{itemize}
\raggedright
  \item \emph{pasi}
  \item it once more, slowly
  \item \emph{pasu}, then \emph{pasi}, and hold the two endings apart
\end{itemize}

\begin{tcolorbox}[breakable,colback=teal!4,colframe=teal!35!black,title={Before you move on}]
What does \ta{பசி} mean? (\textquotedblleft{}hunger\textquotedblright{}.) Say it once more.
\end{tcolorbox}
\section[tākam]{\ta{தாகம்} (tākam) — thirst}

\label{lesson:TA-C61-thirst}



\begin{quote}
Before the new one: what did \emph{pasi} mean?
\end{quote}

\subsection*{You'll want to know: \ta{தாகம்}}
\textbf{\ta{தாகம்}} (\emph{tākam}) — \textquotedblleft{}thirst\textquotedblright{}.

Thirst. Sanskrit-descended, unlike \ta{பசி} beside it, and taken in whole with the \emph{-am} Tamil hangs on such arrivals — the same \emph{-am} already sitting on \ta{புத்தகம்}.

It is answered with \ta{தண்ணீர்} and with nothing else in particular, and a Tamil house offers the \ta{தண்ணீர்} before it is asked for.

\begin{grammarlens}[title={What you've built}]
Two.
\end{grammarlens}

\subsection*{Your turn}
Say these aloud:

\begin{itemize}
\raggedright
  \item \emph{tākam}
  \item it once more, slowly
  \item \emph{pasi}, then \emph{tākam}, and say what each of the two is answered with
\end{itemize}

\begin{tcolorbox}[breakable,colback=teal!4,colframe=teal!35!black,title={Before you move on}]
What does \ta{தாகம்} mean? (\textquotedblleft{}thirst\textquotedblright{}.) Say it once more.
\end{tcolorbox}
\section[tūkkam]{\ta{தூக்கம்} (tūkkam) — sleep}

\label{lesson:TA-C61-sleep}



\begin{quote}
Before the new one: what did \emph{tākam} mean?
\end{quote}

\subsection*{You'll want to know: \ta{தூக்கம்}}
\textbf{\ta{தூக்கம்}} (\emph{tūkkam}) — \textquotedblleft{}sleep\textquotedblright{}.

Sleep as a state — the sleep you get, lose, or need more of. Kannada \ka{ನಿದ್ದೆ} (\emph{nidde}) and Telugu \te{నిద్ర} (\emph{nidra}) both took the Sanskrit word; Tamil kept its own and stayed out of it.

Open it on a long \ta{ஊ}, the same long \ta{ஊ} that opens \ta{ஊசி}, and then hold the \ta{க்க} twice.

\begin{grammarlens}[title={What you've built}]
Three.
\end{grammarlens}

\subsection*{Your turn}
Say these aloud:

\begin{itemize}
\raggedright
  \item \emph{tūkkam}
  \item it once more, slowly
  \item \emph{tūkkam}, then \emph{ūsi}, and hear the same long vowel open both
\end{itemize}

\begin{tcolorbox}[breakable,colback=teal!4,colframe=teal!35!black,title={Before you move on}]
What does \ta{தூக்கம்} mean? (\textquotedblleft{}sleep\textquotedblright{}.) Say it once more.
\end{tcolorbox}
\section[kāyccal]{\ta{காய்ச்சல்} (kāyccal) — a fever}

\label{lesson:TA-C61-fever}



\begin{quote}
Before the new one: what did \emph{tūkkam} mean?
\end{quote}

\subsection*{You'll want to know: \ta{காய்ச்சல்}}
\textbf{\ta{காய்ச்சல்}} (\emph{kāyccal}) — \textquotedblleft{}a fever\textquotedblright{}.

The body running hot. Tamil built the word on its own idea of heating rather than importing one, which is why it looks nothing like Kannada's \ka{ಜ್ವರ} (\emph{jvara}), taken straight from the Sanskrit.

It is the first thing anybody tells a \ta{மருத்துவர்}, and it is said flatly, the way you would report weather.

\begin{grammarlens}[title={What you've built}]
Four.
\end{grammarlens}

\subsection*{Your turn}
Say these aloud:

\begin{itemize}
\raggedright
  \item \emph{kāyccal}
  \item it once more, slowly
  \item \emph{kāyccal}, then \emph{maruttuvar}, and say which one you take to the other
\end{itemize}

\begin{tcolorbox}[breakable,colback=teal!4,colframe=teal!35!black,title={Before you move on}]
What does \ta{காய்ச்சல்} mean? (\textquotedblleft{}a fever\textquotedblright{}.) Say it once more.
\end{tcolorbox}
\section[vali]{\ta{வலி} (vali) — pain}

\label{lesson:TA-C61-pain}



\begin{quote}
Before the new one: what did \emph{kāyccal} mean?
\end{quote}

\subsection*{You'll want to know: \ta{வலி}}
\textbf{\ta{வலி}} (\emph{vali}) — \textquotedblleft{}pain\textquotedblright{}.

Pain, plainly: the ache itself, wherever it sits. Kannada says \ka{ನೋವು} (\emph{nōvu}). Put it with what you already have for the body and it does its work at once — a \ta{வலி} in the \ta{முதுகு}, a \ta{வலி} in the \ta{பல்}.

Keep it away from \emph{vaḻi}, the way or the route, which came up beside \ta{பாதை}. \ta{வலி} takes the plain \ta{ல}; \emph{vaḻi} takes the \ta{ழ} you learned to curl the tongue back for. Two letters English would spell alike, and Tamil holds them apart.

\begin{grammarlens}[title={What you've built}]
Five, and the run is closed: hunger, thirst, sleep, a fever, and pain.
\end{grammarlens}

\subsection*{Your turn}
Say these aloud:

\begin{itemize}
\raggedright
  \item \emph{vali}
  \item it once more, slowly
  \item all five in order — \emph{pasi}, \emph{tākam}, \emph{tūkkam}, \emph{kāyccal}, \emph{vali}
\end{itemize}

\begin{tcolorbox}[breakable,colback=teal!4,colframe=teal!35!black,title={Before you move on}]
What does \ta{வலி} mean? (\textquotedblleft{}pain\textquotedblright{}.) Say it once more.
\end{tcolorbox}
